\documentclass[12pt,a4paper]{article}
\usepackage[margin=25mm, headheight=15pt, headsep=10pt]{geometry}
\usepackage{fontspec}
\usepackage[table]{xcolor} % Note: table option is required for rowcolors
\usepackage{titlesec}
\usepackage{tocloft}
\usepackage{fancyhdr}
\usepackage{booktabs}
\usepackage{tcolorbox}
\tcbuselibrary{skins, breakable}
\usepackage[colorlinks=true, linkcolor=emerald, urlcolor=emerald, bookmarks=true]{hyperref}

\definecolor{emerald}{RGB}{0,102,51}
\definecolor{darkred}{RGB}{139,0,0}
\definecolor{lightgray}{RGB}{245,245,245}

\usepackage[nil]{babel}
\babelprovide[import, main]{bengali}
\babelprovide[import]{arabic}
\babelfont{rm}[Renderer=HarfBuzz,Scale=1.0]{Noto Serif Bengali}
\babelfont{sf}[Renderer=HarfBuzz]{Noto Sans Bengali}
\babelfont{tt}[Renderer=HarfBuzz]{Noto Sans Bengali}
\babelfont[arabic]{rm}[Renderer=HarfBuzz,Scale=1.1]{Noto Naskh Arabic}

% Section styling
\titleformat{\section}{\Large\bfseries\color{emerald}}{\thesection.}{0.5em}{}
\titleformat{\subsection}{\large\bfseries\color{emerald!85!black}}{\thesubsection.}{0.5em}{}
\titleformat{\subsubsection}{\normalsize\bfseries\color{emerald!70!black}}{\thesubsubsection.}{0.5em}{}
\titlespacing*{\section}{0pt}{18pt}{8pt}

% Fancy Header/Footer setup
\pagestyle{fancy}
\fancyhf{} % clear all header and footer fields
\fancyhead[L]{\color{emerald}\small\textbf{\nouppercase{\leftmark}}}
\fancyfoot[L]{\scriptsize\color{black!50}{\lat{AI}}-সংকলিত, আলেম-যাচাইকৃত নয়}
\fancyfoot[C]{\color{emerald!70!black}\thepage}
\renewcommand{\headrulewidth}{0.4pt}
\renewcommand{\headrule}{\hbox to\headwidth{\color{emerald}\leaders\hrule height \headrulewidth\hfill}}
\renewcommand{\footrulewidth}{0pt}

% Custom boxes
% 1. Ayat/Hadith Box
\newtcolorbox{ayatbox}[1][]{
    enhanced, breakable, 
    colback=emerald!5, 
    colframe=emerald, 
    boxrule=0.5pt, 
    arc=3pt, 
    left=10pt, right=10pt, top=10pt, bottom=10pt,
    #1
}

% 2. Quote/Translation Box
\newtcolorbox{quotebox}[1][]{
    enhanced, breakable,
    frame hidden,
    colback=lightgray,
    borderline west={3pt}{0pt}{emerald},
    left=15pt, right=10pt, top=10pt, bottom=10pt,
    sharp corners,
    #1
}

% 3. AI-generated notice box
\newtcolorbox{warnbox}[1][]{
    enhanced, breakable,
    colback=red!4,
    colframe=darkred,
    boxrule=0.8pt,
    arc=3pt,
    left=10pt, right=10pt, top=10pt, bottom=10pt,
    #1
}

% Commands
\newcommand{\arab}[1]{\foreignlanguage{arabic}{#1}}
\newcommand{\kib}[1]{\textcolor{emerald}{\textbf{#1}}}

% Latin (English) fallback: Noto Serif Bengali has NO Latin glyphs
\newfontfamily\latfont[Renderer=HarfBuzz]{Noto Serif}
\newcommand{\lat}[1]{{\latfont #1}}

\setlength{\parskip}{5pt}
\setlength{\parindent}{0pt}

% Fix for missing bullet in Noto Serif Bengali
\renewcommand{\labelitemi}{$\bullet$}



\begin{document}
\begin{center}{\Large\bfseries\color{emerald} \lat{02}-ছোট-সূরা-সমূহ}\end{center}
\vspace{1em}
\section*{ছোট সূরা সমূহ — শব্দ-শব্দ অর্থ ও ব্যাখ্যা (নামাজে পঠিত সাধারণ সূরা)}

\begin{warnbox}
 \textbf{গুরুত্বপূর্ণ নোটিশ (\lat{Important} \lat{Notice}):} এই কন্টেন্টটি \lat{AI} (কৃত্রিম বুদ্ধিমত্তা) দিয়ে প্রস্তুত ও সংকলিত। \lat{AI} কঠোর সূত্র-মান নীতি মেনে শুধু নির্ভরযোগ্য উৎস থেকে তথ্য সংগ্রহ করেছে — কেবল সহীহ/হাসান হাদিস ও সহীহ সনদের আসার রাখা হয়েছে, দুর্বল সনদ বাদ দেওয়া হয়েছে। তবে এটি কোনো আলেম বা মুহাদ্দিস কর্তৃক যাচাইকৃত নয়।
\end{warnbox}

\medskip\hrule\medskip

\subsection*{১. সূত্র কার্ড}

\begin{table}[h]\centering\small
\begin{tabular}{ll}
বিষয় & বিবরণ \\
\hline
\textbf{পরিসংখ্যান} & ৭টি সূরা — নামাজে প্রথম দুই রাকাতে (বা সব রাকাতে) ফাতিহার পর পঠিত \\
\textbf{উৎস} & সহীহ বুখারি ৭৬২, সহীহ মুসলিম ৪৫১ \\
\textbf{গুরুত্ব} & নামাজের নফল কিরাআত — সূরা ফাতিহার পর পঠন সুন্নাতে মুআক্কাদাহ \\
\hline
\end{tabular}
\end{table}

\medskip\hrule\medskip

\subsection*{২. প্রসঙ্গ ও প্রেক্ষাপট}

নামাজের প্রথম দুই রাকাতে সূরা ফাতিহার পর একটি সূরা পাঠ করা \textbf{সুন্নাতে মুআক্কাদাহ} (অত্যন্ত প্রশংসনীয়)। তৃতীয় ও চতুর্থ রাকাতে কোনো সূরা পড়া \textbf{যায় না} (শাফি'ি মাযহাব অনুযায়ী)। এই সূরাগুলো সংক্ষিপ্ত কিন্তু অর্থবহ — প্রতিটিতে আছে \textbf{\lat{tawheed} (তাওহীদ)}, \textbf{\lat{ikhlas} (ইখলাস)}, \textbf{\lat{akhirah} (আখেরাত)} ও \textbf{\lat{qiyamah} (কিয়ামত)}-র গুরুত্বপূর্ণ শিক্ষা।

\medskip\hrule\medskip

\subsection*{৩. সূরা ১: আল-কাওসার (১০৮)}

\subsubsection*{মূল আরবি}
\begin{center}\arab{إِنَّا أَعْطَيْنَاكَ الْكَوْثَرَ}\end{center}
\begin{center}\arab{فَصَلِّ لِرَبِّكَ وَانْحَرْ}\end{center}
\begin{center}\arab{إِنَّ شَانِئَكَ هُوَ الْأَبْتَرُ}\end{center}

\subsubsection*{বাংলা উচ্চারণ}
\textbf{ইন্না আআতাইনাকাল কাওসার}
\textbf{ফাসাল্লি লিরাব্বিকা ওয়ানহার}
\textbf{ইন্না শানি'আকা হুয়াল আবতার}

\subsubsection*{শব্দ-শব্দ অর্থ}

\begin{table}[h]\centering\small
\begin{tabular}{lll}
আরবি & বাংলা & \lat{English} \\
\hline
\textbf{\arab{إِنَّا}} & নিশ্চয়ই আমরা & \textit{\lat{Indeed}, \lat{We}} \\
\textbf{\arab{أَعْطَيْنَاكَ}} & তোমাকে দিয়েছি & \textit{\lat{have} \lat{given} \lat{you}} \\
\textbf{\arab{الْكَوْثَرَ}} & কাওসার (অপার দান/নেয়ামত/জান্নাতের নদী) & \textit{\lat{al-Kawthar} (\lat{abundance} \lat{of} \lat{goodness})} \\
\textbf{\arab{فَصَلِّ}} & তাই নামাজ পড় & \textit{\lat{so} \lat{pray}} \\
\textbf{\arab{لِرَبِّكَ}} & তোমার রবের জন্য & \textit{\lat{for} \lat{your} \lat{Lord}} \\
\textbf{\arab{وَانْحَرْ}} & ও কুরবানি কর (কিংবা মাথা নত কর) & \textit{\lat{and} \lat{sacrifice} (\lat{bow} \lat{down})} \\
\textbf{\arab{إِنَّ}} & নিশ্চয়ই & \textit{\lat{Indeed}} \\
\textbf{\arab{شَانِئَكَ}} & তোমার শত্রু & \textit{\lat{your} \lat{enemy}} \\
\textbf{\arab{هُوَ}} & সেই & \textit{\lat{he}} \\
\textbf{\arab{الْأَبْتَرُ}} & বংশহীন/কুফরির কলঙ্কে ক্ষুণ্ণ & \textit{\lat{cut} \lat{off} \lat{from} \lat{posterity}/\lat{effaced}} \\
\hline
\end{tabular}
\end{table}

\subsubsection*{পূর্ণ বাংলা অনুবাদ}
\textbf{\lat{“}নিশ্চয়ই আমি তোমাকে (হে মুহাম্মদ \arab{ﷺ}) কাওসার দিয়েছি। তাই তোমার রবের জন্য নামাজ পড় এবং কুরবানি কর। নিশ্চয়ই তোমার শত্রুই বংশহীন।\lat{”}}

\subsubsection*{সংক্ষিপ্ত ব্যাখ্যা}
\textbf{\lat{al-Kawthar} (কাওসার)} হলো অপার দান — জান্নাতের নদী, আমলে সালেহ, সন্তানসন্ততি, ও উম্মাহর জন্য রহমাত। পূর্ববর্তী আরবের শত্রুরা মুহাম্মদ \arab{ﷺ}-কে \textbf{\lat{abtar} (বংশহীন)} বলে অপমান করত — কারণ তাঁর প্রথম সন্তানরা শৈশবেই মারা গিয়েছিল। কিন্তু আল্লাহ প্রতিশ্রুতি দিলেন: তোমার শত্রুই আসলে \textbf{\lat{abtar} (বংশহীন)} — তাদের কুফরির কোনো চূড়ান্ত ফল নেই।

\medskip\hrule\medskip

\subsection*{৪. সূরা ২: আল-আসর (১০৩)}

\subsubsection*{মূল আরবি}
\begin{center}\arab{وَالْعَصْرِ}\end{center}
\begin{center}\arab{إِنَّ الْإِنْسَانَ لَفِي خُسْرٍ}\end{center}
\begin{center}\arab{إِلَّا الَّذِينَ آمَنُوا وَعَمِلُوا الصَّالِحَاتِ وَتَوَاصَوْا بِالْحَقِّ وَتَوَاصَوْا بِالصَّبْرِ}\end{center}

\subsubsection*{বাংলা উচ্চারণ}
\textbf{ওয়াল আসর}
\textbf{ইন্নাল ইনসানা লাফী খুসর}
\textbf{ইল্লাল্লাজীনা আমানু ওয়া আমিলুস সালিহাতি ওয়া তাওয়াসাউ বিল হাক্বকী ওয়া তাওয়াসাউ বিস সাবর}

\subsubsection*{শব্দ-শব্দ অর্থ}

\begin{table}[h]\centering\small
\begin{tabular}{lll}
আরবি & বাংলা & \lat{English} \\
\hline
\textbf{\arab{وَالْعَصْرِ}} & সময়ের শপথ & \textit{\lat{By} \lat{the} \lat{time}} \\
\textbf{\arab{إِنَّ}} & নিশ্চয়ই & \textit{\lat{Indeed}} \\
\textbf{\arab{الْإِنْسَانَ}} & মানুষ & \textit{\lat{man}/\lat{humanity}} \\
\textbf{\arab{لَفِي}} & নিশ্চয়ই ... আছে & \textit{\lat{is}} \\
\textbf{\arab{خُسْرٍ}} & ক্ষতিতে/লোকসানে & \textit{\lat{in} \lat{loss}} \\
\textbf{\arab{إِلَّا}} & কিন্তু & \textit{\lat{except}} \\
\textbf{\arab{الَّذِينَ}} & যারা & \textit{\lat{those} \lat{who}} \\
\textbf{\arab{آمَنُوا}} & ঈমান এনেছে & \textit{\lat{believed}} \\
\textbf{\arab{وَعَمِلُوا}} & ও আমল করেছে & \textit{\lat{and} \lat{did}} \\
\textbf{\arab{الصَّالِحَاتِ}} & সৎকর্ম & \textit{\lat{righteous} \lat{deeds}} \\
\textbf{\arab{وَتَوَاصَوْا}} & ও পরস্পর উপদেশ দিয়েছে & \textit{\lat{and} \lat{advised} \lat{each} \lat{other}} \\
\textbf{\arab{بِالْحَقِّ}} & সত্য দিয়ে & \textit{\lat{to} \lat{truth}} \\
\textbf{\arab{وَتَوَاصَوْا}} & ও পরস্পর উপদেশ দিয়েছে & \textit{\lat{and} \lat{advised} \lat{each} \lat{other}} \\
\textbf{\arab{بِالصَّبْرِ}} & ধৈর্য দিয়ে & \textit{\lat{to} \lat{patience}} \\
\hline
\end{tabular}
\end{table}

\subsubsection*{পূর্ণ বাংলা অনুবাদ}
\textbf{\lat{“}সময়ের শপথ! নিশ্চয়ই মানুষ ক্ষতিতেই আছে। তবে যারা ঈমান এনেছে, সৎকর্ম করেছে, সত্যের পরামর্শ দিয়েছে ও ধৈর্যের পরামর্শ দিয়েছে — তারা ব্যতীত।\lat{”}}

\subsubsection*{সংক্ষিপ্ত ব্যাখ্যা}
মোট ৩ আয়াত, কিন্তু সমগ্র কুরআনের সারকে ৪টি \textbf{\lat{key} \lat{principles} (মূলনীতি)}-তে ধরে ফেলেছে: (১) \textbf{\lat{Iman} (ঈমান)} — বিশ্বাস, (২) \textbf{\lat{Amal-e-Salih} (আমলে সালেহ)} — সৎকর্ম, (৩) \textbf{\lat{Tawasi} \lat{bil-Haq} (তাওয়াসী বিল হাক্ব)} — সত্য শেখানো, (৪) \textbf{\lat{Tawasi} \lat{bil-Sabr} (তাওয়াসী বিস সাবর)} — ধৈর্য শেখানো। হাদিসে এসেছে: \textbf{\lat{“}যে ব্যক্তি সূরা আসর পড়ে (বুঝে) পরে, তার জন্য জান্নাতের প্রতিশ্রুতি আছে।\lat{”}} (সুনানে আবু দাউদ ৫০০৫, তিরমিযী ৩০২৯ — সহীহ)।

\medskip\hrule\medskip

\subsection*{৫. সূরা ৩: আল-ইখলাস (১১২)}

\subsubsection*{মূল আরবি}
\begin{center}\arab{قُلْ هُوَ اللَّهُ أَحَدٌ}\end{center}
\begin{center}\arab{اللَّهُ الصَّمَدُ}\end{center}
\begin{center}\arab{لَمْ يَلِدْ وَلَمْ يُولَدْ}\end{center}
\begin{center}\arab{وَلَمْ يَكُنْ لَهُ كُفُوًا أَحَدٌ}\end{center}

\subsubsection*{বাংলা উচ্চারণ}
\textbf{কুল হুয়াল্লাহু আহাদ}
\textbf{আল্লাহুস সামাদ}
\textbf{লাম ইয়ালিদ ওয়ালাম ইউলাদ}
\textbf{ওয়ালাম ইয়াকুন লাহু কুফুয়ান আহাদ}

\subsubsection*{শব্দ-শব্দ অর্থ}

\begin{table}[h]\centering\small
\begin{tabular}{lll}
আরবি & বাংলা & \lat{English} \\
\hline
\textbf{\arab{قُلْ}} & বলো & \textit{\lat{Say}} \\
\textbf{\arab{هُوَ}} & তিনি & \textit{\lat{He}} \\
\textbf{\arab{اللَّهُ}} & আল্লাহ & \textit{\lat{Allah}} \\
\textbf{\arab{أَحَدٌ}} & একক & \textit{\lat{is} \lat{One}} \\
\textbf{\arab{اللَّهُ}} & আল্লাহ & \textit{\lat{Allah}} \\
\textbf{\arab{الصَّمَدُ}} & সমস্ত জগত তাঁর উপর নির্ভরশীল, তিনি কাউকের উপর নির্ভরশীল নন & \textit{\lat{the} \lat{Eternal}, \lat{Self-sufficient}} \\
\textbf{\arab{لَمْ}} & না & \textit{\lat{He} \lat{neither}} \\
\textbf{\arab{يَلِدْ}} & জন্ম দিয়েছেন & \textit{\lat{begets}} \\
\textbf{\arab{وَلَمْ}} & ও না & \textit{\lat{nor}} \\
\textbf{\arab{يُولَدْ}} & জন্মগ্রহণ করেছেন & \textit{\lat{is} \lat{born}} \\
\textbf{\arab{وَلَمْ}} & ও না & \textit{\lat{nor} \lat{is} \lat{there}} \\
\textbf{\arab{يَكُنْ}} & থাকে & \textit{\lat{any}} \\
\textbf{\arab{لَهُ}} & তাঁর জন্য & \textit{\lat{for} \lat{Him}} \\
\textbf{\arab{كُفُوًا}} & সমতুল্য & \textit{\lat{equal}} \\
\textbf{\arab{أَحَدٌ}} & একজনও & \textit{\lat{anyone}} \\
\hline
\end{tabular}
\end{table}

\subsubsection*{পূর্ণ বাংলা অনুবাদ}
\textbf{\lat{“}বলো, তিনি আল্লাহ, একক। আল্লাহ সমদ (সমস্ত জগত তাঁর উপর নির্ভরশীল)। তিনি কাউকে জন্ম দেননি এবং তিনি জন্মগ্রহণও করেননি। আর তাঁর সমতুল্য কেউ নেই।\lat{”}}

\subsubsection*{সংক্ষিপ্ত ব্যাখ্যা}
সূরা ইখলাস \textbf{\lat{tawheed} (তাওহীদ)}-এর মূল ঘোষণা — আল্লাহ একক, কোনো সন্তান নেই, কোনো পিতা নেই, কোনো সমতুল্য নেই। হাদিসে এসেছে: \textbf{\lat{“}যে ব্যক্তি একদিনে তিনবার সূরা ইখলাস পড়ে, সে যেন চারিদিকের সকল মানুষের দোয়ার সমান পুণ্য পেয়েছে।\lat{”}} (সুনানে আবু দাউদ ১৪৬৮, তিরমিযী ৩৫৭৮ — সহীহ)।

\medskip\hrule\medskip

\subsection*{৬. সূরা ৪: আল-ফালাক (১১৩)}

\subsubsection*{মূল আরবি}
\begin{center}\arab{قُلْ أَعُوذُ بِرَبِّ الْفَلَقِ}\end{center}
\begin{center}\arab{مِن شَرِّ مَا خَلَقَ}\end{center}
\begin{center}\arab{وَمِن شَرِّ غَاسِقٍ إِذَا وَقَبَ}\end{center}
\begin{center}\arab{وَمِن شَرِّ النَّفَّاثَاتِ فِي الْعُقَدِ}\end{center}
\begin{center}\arab{وَمِن شَرِّ حَاسِدٍ إِذَا حَسَدَ}\end{center}

\subsubsection*{বাংলা উচ্চারণ}
\textbf{কুল আউযু বিরাব্বিল ফালাক}
\textbf{মিন শাররি মা খালাক}
\textbf{ওয়া মিন শাররি গাসিকিইন ইযা ওয়াকাব}
\textbf{ওয়া মিন শাররিন নাফফাসাতি ফিল উকাদ}
\textbf{ওয়া মিন শাররি হাসিদিইন ইযা হাসাদ}

\subsubsection*{শব্দ-শব্দ অর্থ}

\begin{table}[h]\centering\small
\begin{tabular}{lll}
আরবি & বাংলা & \lat{English} \\
\hline
\textbf{\arab{قُلْ}} & বলো & \textit{\lat{Say}} \\
\textbf{\arab{أَعُوذُ}} & আমি আশ্রয় চাই & \textit{\lat{I} \lat{seek} \lat{refuge}} \\
\textbf{\arab{بِرَبِّ}} & পালনকর্তার কাছে & \textit{\lat{in} \lat{the} \lat{Lord}} \\
\textbf{\arab{الْفَلَقِ}} & ভোরের (সকালের আলো) & \textit{\lat{of} \lat{daybreak}} \\
\textbf{\arab{مِن} \arab{شَرِّ}} & মন্দ থেকে & \textit{\lat{from} \lat{the} \lat{evil}} \\
\textbf{\arab{مَا} \arab{خَلَقَ}} & যা তিনি সৃষ্টি করেছেন & \textit{\lat{of} \lat{that} \lat{He} \lat{has} \lat{created}} \\
\textbf{\arab{وَمِن} \arab{شَرِّ}} & ও মন্দ থেকে & \textit{\lat{and} \lat{from} \lat{the} \lat{evil}} \\
\textbf{\arab{غَاسِقٍ}} & অন্ধকারের (রাতের) & \textit{\lat{of} \lat{darkness} \lat{when} \lat{it} \lat{settles}} \\
\textbf{\arab{إِذَا} \arab{وَقَبَ}} & যখন তা ঘন হয় & \textit{\lat{when} \lat{it} \lat{overtakes}} \\
\textbf{\arab{وَمِن} \arab{شَرِّ}} & ও মন্দ থেকে & \textit{\lat{and} \lat{from} \lat{the} \lat{evil}} \\
\textbf{\arab{النَّفَّاثَاتِ}} & নুনছানোদের (মন্ত্র পাঠকারিণীদের) & \textit{\lat{of} \lat{those} \lat{who} \lat{blow}} \\
\textbf{\arab{فِي} \arab{الْعُقَ}দ\arab{ِ}} & গিঁটে (পরিবেষ্টিত) & \textit{\lat{on} \lat{knots}} \\
\textbf{\arab{وَمِن} \arab{شَرِّ}} & ও মন্দ থেকে & \textit{\lat{and} \lat{from} \lat{the} \lat{evil}} \\
\textbf{\arab{حَاسِدٍ}} & হিংসুকের & \textit{\lat{of} \lat{an} \lat{envier}} \\
\textbf{\arab{إِذَا} \arab{حَسَدَ}} & যখন সে হিংসা করে & \textit{\lat{when} \lat{he} \lat{envies}} \\
\hline
\end{tabular}
\end{table}

\subsubsection*{পূর্ণ বাংলা অনুবাদ}
\textbf{\lat{“}বলো, আমি প্রভাত-কালের (ভোরের আলোর) রবের কাছে আশ্রয় চাই — সৃষ্ট সকল কিছুর অনিষ্ট থেকে, এবং ঘন অন্ধকারের (রাতের) অনিষ্ট থেকে, এবং গিঁটে মন্ত্র পাঠকারিণীদের অনিষ্ট থেকে, এবং হিংসুকের অনিষ্ট থেকে — যখন সে হিংসা করে।\lat{”}}

\subsubsection*{সংক্ষিপ্ত ব্যাখ্যা}
সূরা ফালাক ও সূরা নাস মিলে \textbf{মুযাওযীযাতান (দুটি সুরক্ষা সূরা)}। সূরা ফালাকে আছে \textbf{সকল অনিষ্ট থেকে সুরক্ষা} — প্রকাশ্য ও গোপন শর, জাদু, হিংসা। সূরা নাসে আছে \textbf{কুমন্ত্রণা ও ভ্রান্তি থেকে সুরক্ষা}। হাদিসে এসেছে: \textbf{\lat{“}দুটি সূরা পড়ো — ফালাক ও নাস — প্রতিটি রাকাতের শেষে, আল্লাহর শরণাপন্ন হয়ে।\lat{”}} (সুনানে আবু দাউদ ৫০০২, নাসায়ী ১৩৩৭ — সহীহ)।

\medskip\hrule\medskip

\subsection*{৭. সূরা ৫: আন-নাস (১১৪)}

\subsubsection*{মূল আরবি}
\begin{center}\arab{قُلْ أَعُوذُ بِرَبِّ النَّاسِ}\end{center}
\begin{center}\arab{مَلِكِ النَّاسِ}\end{center}
\begin{center}\arab{إِلَهِ النَّاسِ}\end{center}
\begin{center}\arab{مِن شَرِّ الْوَسْوَاسِ الْخَنَّاسِ}\end{center}
\begin{center}\arab{الَّذِي يُوَسْوِسُ فِي صُدُورِ النَّاسِ}\end{center}
\begin{center}\arab{مِنَ الْجِنَّةِ وَالنَّاسِ}\end{center}

\subsubsection*{বাংলা উচ্চারণ}
\textbf{কুল আউযু বিরাব্বিন নাস}
\textbf{মালিকিন নাস}
\textbf{ইলাহিন নাস}
\textbf{মিন শাররিল ওয়াসওয়াসিল খান্নাস}
\textbf{আল্লাজী ইউওয়াসওয়িসু ফী সুদুরিন নাস}
\textbf{মিনাল জিন্নাতি ওয়ান নাস}

\subsubsection*{শব্দ-শব্দ অর্থ}

\begin{table}[h]\centering\small
\begin{tabular}{lll}
আরবি & বাংলা & \lat{English} \\
\hline
\textbf{\arab{قُلْ}} & বলো & \textit{\lat{Say}} \\
\textbf{\arab{أَعُوذُ}} & আমি আশ্রয় চাই & \textit{\lat{I} \lat{seek} \lat{refuge}} \\
\textbf{\arab{بِرَبِّ}} & পালনকর্তার কাছে & \textit{\lat{in} \lat{the} \lat{Lord}} \\
\textbf{\arab{النَّاسِ}} & মানুষদের & \textit{\lat{of} \lat{mankind}} \\
\textbf{\arab{مَلِكِ}} & মালিক & \textit{\lat{the} \lat{King}} \\
\textbf{\arab{إِلَهِ}} & ইলাহ & \textit{\lat{the} \lat{God}} \\
\textbf{\arab{مِن} \arab{شَرِّ}} & মন্দ থেকে & \textit{\lat{from} \lat{the} \lat{evil}} \\
\textbf{\arab{الْوَسْوَاسِ}} & ফিসফিসানো/কুমন্ত্রণা দানকারীর & \textit{\lat{of} \lat{the} \lat{whisperer}} \\
\textbf{\arab{الْخَنَّاسِ}} & গায়েব হয়ে যাওয়াকারীর (লুকিয়ে কুমন্ত্রণা দেয় এমন) & \textit{\lat{the} \lat{retreating} \lat{whisperer}} \\
\textbf{\arab{الَّذِي}} & যে & \textit{\lat{who}} \\
\textbf{\arab{يُوَسْوِسُ}} & কুমন্ত্রণা দেয় & \textit{\lat{whispers}} \\
\textbf{\arab{فِي}} & ...র ভেতরে & \textit{\lat{into}} \\
\textbf{\arab{صُدُورِ}} & বুকের & \textit{\lat{the} \lat{breasts}} \\
\textbf{\arab{مِنَ} \arab{الْجِنَّةِ}} & জিন থেকে & \textit{\lat{from} \lat{among} \lat{the} \lat{jinn}} \\
\textbf{\arab{وَالنَّاسِ}} & ও মানুষ থেকে & \textit{\lat{and} \lat{mankind}} \\
\hline
\end{tabular}
\end{table}

\subsubsection*{পূর্ণ বাংলা অনুবাদ}
\textbf{\lat{“}বলো, আমি মানুষের পালনকর্তার — মানুষের মালিকের — মানুষের ইলাহের — কাছে আশ্রয় চাই — ফিসফিসানো/কুমন্ত্রণা দানকারীর অনিষ্ট থেকে, যে মানুষের বুকে কুমন্ত্রণা দেয় — জিন ও মানুষের মধ্য থেকে।\lat{”}}

\subsubsection*{সংক্ষিপ্ত ব্যাখ্যা}
সূরা নাসে আছে \textbf{কুমন্ত্রণা (\lat{waswasah})} থেকে সুরক্ষা — যা শয়তান, জিন ও মানুষ থেকে আসে। "ইলাহ" মানে ইবাদতের যোগ্য — এই শব্দে আছে \textbf{\lat{divine} \lat{protection} (আল্লাহর সুরক্ষা)} ও \textbf{\lat{trust} (ভরসা)}। সূরা ফালাক ও নাস মিলে নামাজের শেষে পড়া হয় (সামানি রায়ি), বাম দিকে ৩ হারফাহ বিসমিল্লাহ পড়ে শয়ান করে।

\medskip\hrule\medskip

\subsection*{৮. সূরা ৬: আল-কুরাইশ (১০৬)}

\subsubsection*{মূল আরবি}
\begin{center}\arab{لِإِيلَافِ قُرَيْشٍ}\end{center}
\begin{center}\arab{إِيلَافِهِمْ رِحْلَةَ الشِّتَاءِ وَالصَّيْفِ}\end{center}
\begin{center}\arab{فَلْيَعْبُدُوا رَبَّ هَذَا الْبَيْتِ}\end{center}
\begin{center}\arab{الَّذِي أَطْعَمَهُمْ مِنْ جُوعٍ وَآمَنَهُمْ مِنْ خَوْفٍ}\end{center}

\subsubsection*{বাংলা উচ্চারণ}
\textbf{লিইলাফি কুরাইশ}
\textbf{ইলাফিহিম রিহলাতাশ শিতায়ি ওয়াস সাইফ}
\textbf{ফালইয়াবুদু রাব্বা হাযাল বাইত}
\textbf{আল্লাজী আতআমাহুম মিন জুউন ওয়া আযামাহুম মিন খাউফ}

\subsubsection*{শব্দ-শব্দ অর্থ}

\begin{table}[h]\centering\small
\begin{tabular}{lll}
আরবি & বাংলা & \lat{English} \\
\hline
\textbf{\arab{لِإِيلَافِ}} & অভ্যাসের জন্য & \textit{\lat{For} \lat{the} \lat{accustomed}} \\
\textbf{\arab{قُرَيْشٍ}} & কুরাইশদের & \textit{\lat{of} \lat{Quraysh}} \\
\textbf{\arab{إِيلَافِهِمْ}} & তাদের অভ্যাস & \textit{\lat{their} \lat{accustomed}} \\
\textbf{\arab{رِحْلَةَ}} & যাত্রা & \textit{\lat{caravan}} \\
\textbf{\arab{الشِّتَاءِ}} & শীতকালের & \textit{\lat{of} \lat{winter}} \\
\textbf{\arab{وَالصَّيْفِ}} & ও গ্রীষ্মকালের & \textit{\lat{and} \lat{summer}} \\
\textbf{\arab{فَلْيَعْبُدُوا}} & তাই তারা যেন ইবাদত করে & \textit{\lat{so} \lat{let} \lat{them} \lat{worship}} \\
\textbf{\arab{رَبَّ}} & পালনকর্তার & \textit{\lat{the} \lat{Lord}} \\
\textbf{\arab{هَذَا} \arab{الْبَيْتِ}} & এই ঘরের (কাবা) & \textit{\lat{of} \lat{this} \lat{House}} \\
\textbf{\arab{الَّذِي}} & যিনি & \textit{\lat{Who}} \\
\textbf{\arab{أَطْعَمَهُمْ}} & তাদের খাওয়ালেন & \textit{\lat{fed} \lat{them}} \\
\textbf{\arab{مِنْ} \arab{جُوعٍ}} & ক্ষুধা থেকে & \textit{\lat{from} \lat{hunger}} \\
\textbf{\arab{وَآمَنَهُمْ}} & ও তাদের নিরাপদ করলেন & \textit{\lat{and} \lat{gave} \lat{them} \lat{security}} \\
\textbf{\arab{مِنْ} \arab{خَوْفٍ}} & ভয় থেকে & \textit{\lat{from} \lat{fear}} \\
\hline
\end{tabular}
\end{table}

\subsubsection*{পূর্ণ বাংলা অনুবাদ}
\textbf{\lat{“}কুরাইশদের অভ্যাস (স্থিরতা) — শীত ও গ্রীষ্মের যাত্রায় তাদের অভ্যাস — তাই তারা যেন এই ঘরের (কাবা) পালনকর্তার ইবাদত করে — যিনি তাদের ক্ষুধা থেকে খাওয়ালেন এবং ভয় থেকে নিরাপত্তা দিলেন।\lat{”}}

\subsubsection*{সংক্ষিপ্ত ব্যাখ্যা}
কুরাইশরা কাবার পাশে বসবাস করত এবং শীত-গ্রীষ্মে বাণিজ্যযাত্রায় বের হত। আল্লাহ তাদের সমৃদ্ধি দিলেন — তাই তারা \textbf{\lat{gratitude} (কৃতজ্ঞতা)} প্রকাশ করুক। এই সূরায় আছে \textbf{\lat{favor} (অনুগ্রহ)} স্মরণ ও \textbf{\lat{worship} (ইবাদত)}-র আদেশ।

\medskip\hrule\medskip

\subsection*{৯. সূরা ৭: আল-মাউন (১০৭)}

\subsubsection*{মূল আরবি}
\begin{center}\arab{أَرَأَيْتَ الَّذِي يُكَذِّبُ بِالدِّينِ}\end{center}
\begin{center}\arab{فَذَلِكَ الَّذِي يَدُعُّ الْيَتِيمَ}\end{center}
\begin{center}\arab{وَلَا يَحُضُّ عَلَىٰ طَعَامِ الْمِسْكِينِ}\end{center}
\begin{center}\arab{فَوَيْلٌ لِلْمُصَلِّينَ}\end{center}
\begin{center}\arab{الَّذِينَ هُمْ عَنْ صَلَاتِهِمْ سَاهُونَ}\end{center}
\begin{center}\arab{الَّذِينَ هُمْ يُرَاءُونَ}\end{center}
\begin{center}\arab{وَيَمْنَعُونَ الْمَاعُونَ}\end{center}

\subsubsection*{বাংলা উচ্চারণ}
\textbf{আরাইতাল্লাজী ইউকাজ্জিবু বিদ দীন}
\textbf{ফাযালিকাল্লাজী ইয়াদু' উল ইয়াতীম}
\textbf{ওয়ালা ইয়াহুদ্দু আলা তা'আমিল মিসকীন}
\textbf{ফাওয়াইলুন লিল মুসাল্লীন}
\textbf{আল্লাজীনা হুম আন সালাতিহিম সাহুন}
\textbf{আল্লাজীনা হুম ইউরাউন}
\textbf{ওয়াইয়ামনা' উনাল মা'উন}

\subsubsection*{শব্দ-শব্দ অর্থ}

\begin{table}[h]\centering\small
\begin{tabular}{lll}
আরবি & বাংলা & \lat{English} \\
\hline
\textbf{\arab{أَرَأَيْتَ}} & তুমি কি দেখেছ & \textit{\lat{Have} \lat{you} \lat{seen}} \\
\textbf{\arab{الَّذِي}} & সেই ব্যক্তি & \textit{\lat{the} \lat{one} \lat{who}} \\
\textbf{\arab{يُكَذِّبُ}} & মিথ্যা বলে & \textit{\lat{denies}} \\
\textbf{\arab{بِالدِّينِ}} & ধর্মের প্রতি & \textit{\lat{the} \lat{religion}} \\
\textbf{\arab{فَذَلِكَ}} & তাই সেই & \textit{\lat{That} \lat{is} \lat{the} \lat{one}} \\
\textbf{\arab{يَدُعُّ}} & ধাক্কা দেয়/তাড়ায় & \textit{\lat{who} \lat{pushes} \lat{away}} \\
\textbf{\arab{الْيَتِ}ম\arab{َ}} & এতীমকে & \textit{\lat{the} \lat{orphan}} \\
\textbf{\arab{وَلَا}} & ও না & \textit{\lat{and} \lat{does} \lat{not}} \\
\textbf{\arab{يَحُضُّ}} & উৎসাহিত করে & \textit{\lat{encourage}} \\
\textbf{\arab{عَلَىٰ}} & ...র প্রতি & \textit{\lat{for}} \\
\textbf{\arab{طَعَامِ}} & খাদ্য & \textit{\lat{the} \lat{feeding} \lat{of}} \\
\textbf{\arab{الْمِسْكِينِ}} & দুঃখী/গরিবকে & \textit{\lat{the} \lat{poor}} \\
\textbf{\arab{فَوَيْلٌ}} & তাই ধ্বংস/বিপদ & \textit{\lat{So} \lat{woe}} \\
\textbf{\arab{لِلْمُصَلِّينَ}} & নামাজীদের জন্য & \textit{\lat{to} \lat{the} \lat{worshippers}} \\
\textbf{\arab{الَّذِينَ} \arab{هُمْ}} & যারা & \textit{\lat{who} \lat{are}} \\
\textbf{\arab{عَنْ} \arab{صَلَاتِهِمْ}} & তাদের নামাজ থেকে & \textit{\lat{unmindful} \lat{of}} \\
\textbf{\arab{سَاهُون}} & সোজা/অবহেলাকারী & \textit{\lat{their} \lat{prayer}} \\
\textbf{\arab{يُرَاءُونَ}} & দেখানোর জন্য নামাজ পড়ে & \textit{\lat{show} \lat{off}} \\
\textbf{\arab{وَيَمْنَعُونَ}} & ও বাধা দেয় & \textit{\lat{and} \lat{withhold}} \\
\textbf{\arab{الْمَاعُونَ}} & সাধারণ সাহায্য/ব্যবহার্য বস্তু & \textit{\lat{the} \lat{small} \lat{kindnesses}} \\
\hline
\end{tabular}
\end{table}

\subsubsection*{পূর্ণ বাংলা অনুবাদ}
\textbf{\lat{“}তুমি কি সেই ব্যক্তির কথা জানো যে ধর্মকে মিথ্যা প্রতিপাদন করে? সেই ব্যক্তিই এতীমকে ধাক্কা দেয় এবং দুঃখীকে খাওয়ানোর জন্য উৎসাহিত করে না। তাই নামাজীদের জন্য ধ্বংস — যারা তাদের নামাজ সম্পর্কে অবহেলাকারী, যারা দেখানোর জন্য নামাজ পড়ে, এবং যারা সাধারণ সাহায্যও বাধা দেয়।\lat{”}}

\subsubsection*{সংক্ষিপ্ত ব্যাখ্যা}
এই সূরায় দুটি \textbf{\lat{key} \lat{lessons} (মূল শিক্ষা)}: (১) \textbf{\lat{Orphan} (এতীম)} ও \textbf{\lat{poor} (মিসকীন)}-দের প্রতি অবহেলা — এটি ধর্মের প্রতি মিথ্যা প্রতিপাদনের লক্ষণ; (২) \textbf{\lat{Prayer} (নামাজ)} — \textbf{\lat{riya} (দেখানো)} বা \textbf{\lat{distraction} (মনোযোগহীনতা)} সহ পড়লে নামাজ গ্রহণযোগ্য হয় না। নামাজে \textbf{\lat{khushu} (খুশু)} থাকা এবং সমাজে \textbf{\lat{help} (সাহায্য)} করা — দুটোই একসাথে যায়।

\medskip\hrule\medskip

\textit{শেষ আপডেট: আগস্ট ২০২৬}

\end{document}
